% !PRINTED 138
\weberpage{138}
%<<CONT
echte Theiler dieser Zahlen sind, schon als metacyklisch erwiesen sind.

Es bleibt nur noch die Untersuchung der Zahlen 60 und 90 übrig. Dass eine einfache Gruppe vom 60\textsuperscript{sten} Grade existirt, wissen wir schon; nämlich die alternirende Gruppe der Permutationen von fünf Buchstaben (Bd.\,I, \S.\,177). Es ist aber noch fraglich, ob dies die einzige ist.

Eine einfache Gruppe $P$ vom Grade 60 enthält nach \S.\,29, I. als Theiler eine Gruppe 5\textsuperscript{ten} Grades und eine Gruppe 3\textsuperscript{ten} Grades. Nimmt man in dem zweiten \gesperrt{Sylow'schen} Satze (\S.\,30) für $Q$ die Gruppe 5\textsuperscript{ten} Grades, so muss $R$ vom 10\textsuperscript{ten} Grade sein, weil der Index von $R$ congruent mit 1 nach dem Modul 5, also nur $= 6$ sein kann. Also kann die Gruppe $P$ vom 60\textsuperscript{sten} Grade nach \S.\,28, 2. als transitive Permutationsgruppe von sechs Ziffern dargestellt werden; sie kann aber keine cyklische Permutation von nur drei Ziffern enthalten, weil sie als einfache Gruppe (nach Bd.\,I, \S.\,158, 2.) primitiv sein muss, und daher, wenn sie einen dreigliedrigen Cyklus enthielte, nach Bd.\,I, \S.\,153, 10. die ganze alternirende Gruppe 360\textsuperscript{sten} Grades enthalten müsste. Wir können also eine der Permutationen 3\textsuperscript{ten} Grades in der Form annehmen:
\[
a = (1, 2, 3) (4, 5, 6).
\]
Hierzu wollen wir eine Permutation 5\textsuperscript{ten} Grades, $b$, fügen, die nur eine cyklische sein kann. Lassen wir darin die Ziffer 6 fehlen, so können wir statt $b$ eine solche Potenz von $b$ nehmen (die ja auch in $P$ vorkommen muss), dass etwa 1, 2 die beiden ersten Ziffern werden, und dann bleiben noch sechs Möglichkeiten für $b$ übrig:
\[
(1, 2, 3, 4, 5), \quad (1, 2, 4, 3, 5), \quad (1, 2, 4, 5, 3),
\]
\[
(1, 2, 3, 5, 4), \quad (1, 2, 5, 3, 4), \quad (1, 2, 5, 4, 3);
\]
von diesen sechs Annahmen sind aber wieder die drei in einer Reihe stehenden nicht wesentlich (d.\,h. nur durch die Bezeichnung) verschieden. Denn ersetzt man z.\,B. bei der zweiten Annahme $b = (1, 2, 4, 3, 5)$ das Element $a$ durch $a^2$ und $b$ durch $b^3$, so erhält man
\[
(1, 3, 2) (4, 6, 5); \quad (1, 3, 2, 5, 4),
\]
was durch Vertauschung der Ziffern 2 mit 3 und 4 mit 5 in die erste Annahme übergeht. Die dritte Annahme $b = (1, 2, 4, 5, 3)$ geht, wenn man $b$ durch $b^4$ und $a$ durch $a^2$ ersetzt und dann
%CONT>>
